### Definitions

For an integer (m\ge 1), let ([m]={1,2,\dots,m}).
Let (K_m^{(3)}) denote the **complete 3-uniform hypergraph** on vertex set ([m]), i.e. the 3-uniform hypergraph with edge set
[
E(K_m^{(3)})=\binom{[m]}{3}.
]
A **2-colouring** of the edges of (K_m^{(3)}) is a function
[
\chi:\binom{[m]}{3}\to{\text{red},\text{blue}}.
]
For an integer (n\ge 1), let (K_n^{(3)}) denote the complete 3-uniform hypergraph on (n) vertices.

A subset (S\subseteq [m]) with (|S|=n) is a **monochromatic copy of (K_n^{(3)})** (under (\chi)) if (\chi) is constant on (\binom{S}{3}); i.e. every triple in (S) is coloured the same.

Define (R_3(n)) to be the least integer (m) such that **every** 2-colouring (\chi:\binom{[m]}{3}\to{\text{red},\text{blue}}) contains a monochromatic copy of (K_n^{(3)}). Equivalently,
[
R_3(n)=\min\Bigl{m\in\mathbb N_{\ge 1}:\ \forall \chi:\binom{[m]}{3}\to{\text{red},\text{blue}}\ \exists S\subseteq[m],,|S|=n,\ \chi\text{ constant on }\binom{S}{3}\Bigr}.
]

### The “literal statement” to be proved/disproved

The prompt asks:

> “Is there some constant (c>0) such that (R_3(n)\ge 2^{2^{cn}})?”

As a declarative claim (the usual way to formalize such a question), the **literal assertion** is:
[
(\star)\qquad \exists c>0\ \forall n\in\mathbb N_{\ge 1},\quad R_3(n)\ \ge\ 2^{,2^{c n}}.
]

### Ambiguity / minimal corrected statement (standard convention)

In the modern literature this is usually posed only for **(n\ge 4)** (since (n\le 3) are trivial edge-cases for 3-uniform cliques), e.g. as an Erdős conjecture. 

A minimal corrected version consistent with that convention is:
[
(\star_{\ge 4})\qquad \exists c>0\ \forall n\in\mathbb N,\ n\ge 4,\quad R_3(n)\ \ge\ 2^{,2^{c n}}.
]

What I will do below: **disprove the literal assertion ((\star))** (as written, with (n\ge 1)). I will also note that ((\star_{\ge 4})) is the intended open problem.

---
